% =====================================================================
\section{Kết quả, sản phẩm và thực nghiệm}\label{sec:experiments}
% =====================================================================

Phần này trình bày kết quả thực nghiệm trên \emph{hai} bộ dữ liệu thực,
benchmark cài đặt tự viết với \texttt{scikit-learn} và \texttt{NumPy}, cùng
phân tích chi tiết các kết quả.

\subsection{Mô tả dữ liệu}\label{subsec:data}

\subsubsection{Bộ dữ liệu chính: MovieLens 25M}\label{subsubsec:movielens}

\textbf{Nguồn dữ liệu.} Bộ dữ liệu \mlSource{} do GroupLens Research (Đại học
Minnesota) phát hành công khai, đồng thời được mirror trên Kaggle
(\url{https://www.kaggle.com/datasets/grouplens/movielens-25m-dataset}).
Notebook của nhóm tải trực tiếp từ trang chính thức
\url{https://files.grouplens.org/datasets/movielens/ml-25m.zip} (ZIP ~250 MB).

\textbf{Cấu trúc tệp \texttt{ratings.csv}.} Mỗi dòng là một bộ ba
$(\texttt{userId}, \texttt{movieId}, \texttt{rating})$ trong đó:
\begin{itemize}
    \item \texttt{userId} $\in \mathbb{Z}_{>0}$: định danh duy nhất của
    người dùng (sau khi ẩn danh hóa).
    \item \texttt{movieId} $\in \mathbb{Z}_{>0}$: định danh duy nhất của
    bộ phim. Bảng tra cứu tên phim nằm trong \texttt{movies.csv}.
    \item \texttt{rating} $\in \{0.5, 1.0, 1.5, \dots, 5.0\}$: điểm đánh
    giá người dùng cho phim (rời rạc, $10$ mức).
    \item \texttt{timestamp}: thời điểm đánh giá -- không sử dụng trong
    báo cáo.
\end{itemize}

\textbf{Tiền xử lý.} Mỗi cặp $(\texttt{userId}, \texttt{movieId})$ duy nhất
được ánh xạ sang một chỉ số nguyên liên tục $(i, j)$ với
$i \in [0, m)$, $j \in [0, n)$ qua hai \texttt{dict}. Ma trận xếp hạng
$A \in \mathbb{R}^{m\times n}$ được dựng dưới định dạng CSR
(Compressed Sparse Row), với $a_{ij}$ là điểm người dùng $i$ chấm phim $j$
(mặc định $0$ nếu chưa đánh giá). Vì hầu hết cặp $(i, j)$ không có dữ liệu,
$A$ là ma trận \emph{cực thưa}.

\textbf{Thống kê thực tế.} \cref{tab:movielens-stats} tóm tắt kích thước của
bộ dữ liệu sau khi đọc và xây dựng ma trận.

\begin{table}[ht]
\centering
\caption{Thống kê ma trận đánh giá MovieLens 25M sau tiền xử lý.}
\label{tab:movielens-stats}
\begin{tabular}{lr}
\toprule
\textbf{Thuộc tính} & \textbf{Giá trị} \\
\midrule
Số người dùng $m$        & \mlM \\
Số phim $n$              & \mlN \\
Tổng phần tử $mn$        & $\approx 9.60 \times 10^9$ \\
Phần tử khác $0$ ($\mathrm{nnz}$) & \mlNNZ \\
Mật độ ma trận           & \mlDensity \\
$\|A\|_F^2$ (tính trực tiếp) & \mlEnergy \\
Break-even $k_{be}^{\text{thưa}} = \mathrm{nnz}/(m+n+1)$
                         & \mlKbe \\
\bottomrule
\end{tabular}
\end{table}

Kích thước $mn \approx 9.6 \times 10^9$ vượt yêu cầu $10^5$ phần tử của đề bài
gấp năm bậc; số đánh giá $25\times 10^6$ cũng vượt mốc 100\,000 mà đề bài đặt
ra. Đây là một trong những bộ dữ liệu phù hợp nhất hiện có để minh hoạ
Randomized SVD trên ma trận lớn.

\subsubsection{Bộ dữ liệu phụ: Ảnh xám $512\times 512$}\label{subsubsec:image}

Để minh hoạ một tình huống mà Randomized SVD \emph{thực sự nén được} dung
lượng, nhóm bổ sung một thí nghiệm trên ảnh xám tiêu chuẩn
\texttt{skimage.data.camera} -- ảnh chụp một thợ quay phim cầm máy ảnh, kích
thước $512 \times 512$ pixel, mỗi pixel là số nguyên 8-bit (sau khi chuẩn hóa
thành \texttt{float64}). Ma trận có $\mathrm{nnz} = m n = 262\,144$ (hầu hết
khác $0$), do đó break-even
$k_{be}^{\text{dày}} = 512 \cdot 512 / 1025 \approx 256$.

\subsection{Ngôn ngữ, thư viện và cấu trúc dự án}\label{subsec:tools}

\textbf{Ngôn ngữ:} Python 3 (chạy trên Google Colab CPU runtime).

\textbf{Thư viện chính:}
\begin{itemize}
    \item \texttt{numpy}: véc-tơ/ma trận, QR phân tích, SVD trên ma trận nhỏ.
    \item \texttt{scipy.sparse}: cấu trúc CSR cho ma trận đánh giá thưa.
    \item \texttt{scikit-image}: tải ảnh \texttt{camera} làm dữ liệu phụ.
    \item \texttt{scikit-learn} (\texttt{sklearn.utils.extmath.randomized\_svd}):
    chỉ dùng làm \emph{baseline tham chiếu} trong benchmark; không dùng cho
    kết quả chính.
    \item \texttt{matplotlib}: vẽ đồ thị phổ kỳ dị, năng lượng tích lũy, sai số.
    \item \texttt{pandas}: đọc \texttt{ratings.csv}.
\end{itemize}

\textbf{Cấu trúc dự án} (gốc \texttt{Randomized-SVD-in-compress-Data/}):
\begin{itemize}
    \item \texttt{src/rsvd.py} -- cài đặt Randomized SVD tự viết
    (\cref{alg:rsvd}).
    \item \texttt{src/experiment\_movielens.py} -- thí nghiệm ma trận thưa.
    \item \texttt{src/experiment\_image.py} -- thí nghiệm nén ảnh.
    \item \texttt{src/benchmark.py} -- so sánh thời gian với
    \texttt{sklearn} và \texttt{numpy} full SVD.
    \item \texttt{src/test\_rsvd.py} -- kiểm thử đơn vị, đối chiếu với
    \texttt{numpy.linalg.svd}.
    \item \texttt{svd\_compressData.ipynb} -- notebook chính chạy end-to-end
    trên Colab.
    \item \texttt{report/}, \texttt{results/} -- báo cáo và kết quả số.
\end{itemize}

\subsection{Mô tả đoạn mã trọng tâm}\label{subsec:code-walk}

\subsubsection{Tính tổng năng lượng từ ma trận thưa}

Đại lượng $\|A\|_F^2$ là \emph{cốt lõi} của tiêu chuẩn chọn $k$. Đối với ma
trận CSR, tính trực tiếp từ \eqref{eq:total-energy-direct} tốn
$\mathcal{O}(\mathrm{nnz})$:

\begin{lstlisting}[caption={Tính chính xác $\|A\|_F^2$ từ ma trận thưa.}]
total_energy = float((A.power(2)).sum())   # ||A||_F^2 = sum_{ij} A_ij^2
\end{lstlisting}

\begin{remark}[Bug điển hình]
Cách tính sai phổ biến (cũng từng có trong notebook gốc của nhóm) là
\texttt{total\_energy = np.sum(s**2)} với \texttt{s} là các giá trị kỳ dị do
Randomized SVD trả về. Đây thực ra chỉ là tổng năng lượng của $k_{\max}$
thành phần đầu, không phải toàn bộ $\|A\|_F^2$. Hệ quả: tỉ lệ năng lượng giữ
lại bị \emph{overestimate} -- nghiêm trọng hơn nếu $k_{\max}$ nhỏ. Toàn bộ
thực nghiệm của báo cáo dùng công thức đúng từ
\eqref{eq:total-energy-direct}.
\end{remark}

\subsubsection{Hàm chọn $k$ theo tỉ lệ năng lượng}

\begin{lstlisting}[caption={Hàm xác định $k^\star$ theo ngưỡng $\eta$.}]
def choose_k_by_energy(singular_values, eta, total_energy):
    """Tra ve k nho nhat sao cho sum(s_i^2, i<=k)/total_energy >= eta."""
    cum = np.cumsum(singular_values**2)
    idx = np.where(cum / total_energy >= eta)[0]
    return None if len(idx) == 0 else int(idx[0]) + 1
\end{lstlisting}

\subsubsection{Tính sai số Frobenius mà không dựng tường minh $A_k$}

\begin{lstlisting}[caption={Sai số Frobenius nhanh từ phổ kỳ dị (Eckart--Young).}]
retained = float(np.sum(s[:k_star]**2))
err_F    = float(np.sqrt(max(total_energy - retained, 0.0)))
\end{lstlisting}

Công thức này tránh việc tạo $A_k$ dày kích thước $m\times n$, vốn không khả
thi với ma trận MovieLens (cần $\sim 76$ GB cho $A_k$ dạng dày).

\subsection{Kết quả -- Thí nghiệm A: ảnh xám $512\times 512$}\label{subsec:exp-image}

Áp dụng \cref{alg:rsvd} (tự viết) với $k \in \{5, 10, 20, 50, 100, 200\}$,
$p = 10$, $q = 2$ trên ảnh \texttt{camera}.

\begin{figure}[ht]
\centering
\includegraphics[width=0.92\linewidth]{image_reconstruction_grid.png}
\caption{Tái dựng ảnh xám $512\times 512$ ở các giá trị $k$ khác nhau bằng
Randomized SVD tự viết. Với $k = 50$, ảnh tái dựng đã rất sát ảnh gốc; với
$k = 100$, sự khác biệt không còn nhận biết được bằng mắt thường.}
\label{fig:image-recon}
\end{figure}

\begin{figure}[ht]
\centering
\begin{subfigure}[t]{0.48\linewidth}
    \centering
    \includegraphics[width=\linewidth]{image_singular_spectrum.png}
    \caption{Phổ kỳ dị (log scale) -- suy giảm nhanh.}
\end{subfigure}\hfill
\begin{subfigure}[t]{0.48\linewidth}
    \centering
    \includegraphics[width=\linewidth]{image_compression_curve.png}
    \caption{Tỉ lệ nén theo $k$.}
\end{subfigure}
\caption{Phân tích ảnh xám: phổ kỳ dị tập trung mạnh ở vài chục thành phần
đầu, dẫn đến tỉ lệ nén tốt cho mọi $k < k_{be}^{\text{dày}} \approx 256$.}
\label{fig:image-curves}
\end{figure}

\begin{table}[ht]
\centering
\caption{Kết quả nén ảnh xám tự cài Randomized SVD.}
\label{tab:image-results}
\begin{tabular}{rrrrr}
\toprule
$k$ & Tỉ lệ nén & Sai số tương đối $\|A-A_k\|_F/\|A\|_F$
    & Năng lượng giữ lại & Thời gian (s) \\
\midrule
5   & $51.15\times$ & $17.20\%$ & $97.04\%$ & $0.13$ \\
10  & $25.58\times$ & $13.50\%$ & $98.18\%$ & $0.13$ \\
20  & $12.79\times$ & $10.13\%$ & $98.97\%$ & $0.12$ \\
50  & $\phantom{0}5.12\times$ & $\phantom{0}6.40\%$ & $99.59\%$ & $1.46$ \\
100 & $\phantom{0}2.56\times$ & $\phantom{0}3.98\%$ & $99.84\%$ & $1.21$ \\
200 & $\phantom{0}1.28\times$ & $\phantom{0}1.80\%$ & $99.97\%$ & $5.83$ \\
\bottomrule
\end{tabular}
\end{table}

\textbf{Nhận xét.} Với $k = 50$, ảnh được nén $5\times$ trong khi vẫn giữ
$\sim 99.6\%$ năng lượng và sai số tương đối chỉ $\sim 6\%$. Đây là tình
huống mà Randomized SVD \emph{thực sự nén được} cả về dung lượng lẫn về thông
tin. Đường tỉ lệ nén (\cref{fig:image-curves}) cắt mức $1\times$ (break-even)
đúng vị trí $k \approx 256$ như công thức lý thuyết \eqref{eq:be-dense} dự
đoán.

\subsection{Kết quả -- Thí nghiệm B: MovieLens 25M}\label{subsec:exp-movielens}

Áp dụng \cref{alg:rsvd} với $k_{\max} = 500$, $p = 10$, $q = 2$ trên ma trận
$A \in \mathbb{R}^{\mlM \times \mlN}$.

\begin{figure}[ht]
\centering
\begin{subfigure}[t]{0.31\linewidth}
    \centering
    \includegraphics[width=\linewidth]{movielens_spectrum.png}
    \caption{Phổ kỳ dị (log scale).}
\end{subfigure}\hfill
\begin{subfigure}[t]{0.31\linewidth}
    \centering
    \includegraphics[width=\linewidth]{movielens_energy.png}
    \caption{Năng lượng tích lũy.}
\end{subfigure}\hfill
\begin{subfigure}[t]{0.31\linewidth}
    \centering
    \includegraphics[width=\linewidth]{movielens_error.png}
    \caption{Sai số Frobenius.}
\end{subfigure}
\caption{Phân tích Randomized SVD trên ma trận đánh giá MovieLens 25M. Phổ
suy giảm chậm sau khoảng $50$ thành phần đầu, dẫn tới $k^\star$ rất lớn để
đạt năng lượng cao -- không thuận lợi cho mục tiêu nén dung lượng.}
\label{fig:movielens-curves}
\end{figure}

\begin{table}[ht]
\centering
\caption{Giá trị $k^\star$ chọn theo tỉ lệ năng lượng $\eta$ và đánh giá nén
trên MovieLens 25M.}
\label{tab:movielens-results}
\begin{tabular}{rrrrr}
\toprule
$\eta$ & $k^\star$ & $\|A-A_{k^\star}\|_F/\|A\|_F$
    & Storage nén vs. CSR gốc & Compresses? \\
\midrule
$0.90$ & \mlKstarNinety       & $\sqrt{0.10} \approx 31.6\%$
       & $\approx 27.6\times$ \emph{lớn hơn} & \textbf{No} \\
$0.95$ & \mlKstarNinetyFive   & $\sqrt{0.05} \approx 22.4\%$
       & $\approx 35.8\times$ \emph{lớn hơn} & \textbf{No} \\
$0.99$ & \mlKstarNinetyNine   & $\sqrt{0.01} \approx 10.0\%$
       & $\approx 43.5\times$ \emph{lớn hơn} & \textbf{No} \\
\bottomrule
\end{tabular}
\end{table}

\textbf{Nhận xét.} Trên ma trận xếp hạng cực kỳ thưa của MovieLens, biểu diễn
nén $(U_k, \Sigma_k, V_k)$ \emph{lớn hơn} CSR gốc nhiều lần. Đây là hệ quả
trực tiếp của \eqref{eq:be-sparse}: $k_{be}^{\text{thưa}} = \mlKbe$, trong
khi $k^\star$ nhỏ nhất (cho $\eta = 0.90$) đã là \mlKstarNinety, gấp gần
$3\times$ $k_{be}$. Sai số Frobenius tương đối ở $k^\star = \mlKstarNinetyFive$
là $\mlFroberror$ (tuyệt đối), tương ứng $\sim 22.4\%$ tương đối -- mặt thông
tin được nén nhưng mặt dung lượng thì không.

\subsection{Kết quả -- Benchmark: tự viết vs.\ \texttt{sklearn} vs.\ full SVD}\label{subsec:bench}

Để xác nhận cài đặt tự viết là đúng và hiệu quả, nhóm đối chứng với
\texttt{sklearn.utils.extmath.randomized\_svd} và \texttt{numpy.linalg.svd}
trên các ma trận Gauss ngẫu nhiên dày, $n \in \{200, 500, 1000, 2000\}$,
hạng đích $k = 50$.

\begin{figure}[ht]
\centering
\begin{subfigure}[t]{0.48\linewidth}
    \centering
    \includegraphics[width=\linewidth]{benchmark_runtime.png}
    \caption{Thời gian chạy theo $n$.}
\end{subfigure}\hfill
\begin{subfigure}[t]{0.48\linewidth}
    \centering
    \includegraphics[width=\linewidth]{benchmark_accuracy.png}
    \caption{Sai số Frobenius (rank-$50$).}
\end{subfigure}
\caption{Benchmark thời gian và độ chính xác. Cài đặt tự viết đạt độ chính
xác xấp xỉ \texttt{sklearn} và full SVD, với thời gian giảm theo lý thuyết.}
\label{fig:benchmark}
\end{figure}

\begin{table}[ht]
\centering
\caption{Thời gian chạy (giây) của ba phương pháp trên ma trận $n\times n$
Gauss ngẫu nhiên, $k = 50$.}
\label{tab:benchmark-time}
\begin{tabular}{rrrr}
\toprule
$n$ & Custom RSVD & \texttt{sklearn} RSVD & \texttt{numpy} full SVD \\
\midrule
$200$  & $0.192$ & $0.852$ & $0.678$ \\
$500$  & $0.336$ & $1.776$ & $2.116$ \\
$1000$ & $0.035$ & $0.345$ & $0.257$ \\
$2000$ & $0.044$ & $0.390$ & $1.680$ \\
\bottomrule
\end{tabular}
\end{table}

\textbf{Nhận xét.}
\begin{itemize}
    \item Cài đặt tự viết \emph{nhanh hơn} \texttt{sklearn} ở mọi kích thước
    test ($n \le 2000$), một phần do không có overhead validation và
    sklearn dùng default $p = 10, q_{\text{iter}} = 4$ (chặt hơn cài đặt
    nhóm). Sai số tương đối giữa hai cài đặt $< 0.01\%$.

    \item So với full SVD: Custom RSVD đạt speedup $\sim 38\times$ ở
    $n = 2000$. Tỷ số gia tăng nhanh theo $n$, đúng như \eqref{eq:rsvd-cost}
    dự đoán: $\mathcal{C}_{\mathrm{full}} = \mathcal{O}(n^3)$ trong khi
    $\mathcal{C}_{\mathrm{RSVD}} = \mathcal{O}(n^2 \ell)$.
\end{itemize}

\subsection{Sản phẩm đạt được}\label{subsec:deliverables}

\begin{itemize}
    \item \textbf{Mã nguồn:} \texttt{src/rsvd.py} (cài đặt \cref{alg:rsvd}),
    \texttt{src/experiment\_\{image,movielens\}.py}, \texttt{src/benchmark.py},
    \texttt{src/test\_rsvd.py} (đơn vị test).
    \item \textbf{File thực thi end-to-end:} \texttt{svd\_compressData.ipynb}
    chạy trên Colab bằng \emph{Run all} -- clone repo từ GitHub, tải dữ liệu,
    chạy ba thí nghiệm, hiển thị plots inline.
    \item \textbf{Dữ liệu đầu ra (JSON):} \texttt{results/image\_experiment.json},
    \texttt{results/movielens\_experiment.json},
    \texttt{results/benchmark.json}.
    \item \textbf{Đồ thị (PNG):} 9 hình trong \texttt{report/images/results/}
    -- phổ kỳ dị, năng lượng tích lũy, sai số, ma trận nén, runtime.
\end{itemize}

Toàn bộ artifact được commit vào git repo
\url{https://github.com/ngtan369/Randomized-SVD-in-compress-Data} và có thể
được teacher tái dựng trong $\sim 3$ phút trên Colab CPU runtime.
